\documentclass[a4paper,12pt]{article}

\usepackage[utf8]{inputenc}
\usepackage[russian]{babel}
\usepackage{graphicx}
\usepackage{listings}
\usepackage{amsmath}

\title{Лабораторная работа \\ Перемножение квадратных матриц}
\author{Никишин Владислав \\ группа 6312-100503D}
\date{}

\begin{document}

\maketitle

\section{Постановка задачи}

Необходимо разработать программу на языке C/C++ для перемножения двух квадратных матриц.

Исходные данные: файлы содержащие значения матриц.

Выходные данные:
\begin{itemize}
\item результирующая матрица
\item время выполнения программы
\item объём задачи
\end{itemize}

Также требуется выполнить автоматизированную проверку результатов вычислений с помощью стороннего ПО.

\section{Описание алгоритма}

Перемножение матриц выполняется по формуле:

\[
C_{ij} = \sum_{k=1}^{n} A_{ik} B_{kj}
\]

где:

\begin{itemize}
\item A — первая матрица
\item B — вторая матрица
\item C — результирующая матрица
\end{itemize}

Алгоритм реализован с использованием трёх вложенных циклов.

\section{Исходный код программы}

Основная программа реализована на языке C++.

\begin{lstlisting}[language=C++]
for(int i=0;i<n;i++)
    for(int j=0;j<n;j++)
        for(int k=0;k<n;k++)
            C[i][j] += A[i][k] * B[k][j];
\end{lstlisting}

\section{Экспериментальные результаты}

Проведены эксперименты для различных размеров матриц.

\begin{center}
\begin{tabular}{|c|c|}
\hline
Размер матрицы & Время выполнения \\
\hline
100x100 & 0.01 s \\
300x300 & 0.14 s \\
500x500 & 0.52 s \\
\hline
\end{tabular}
\end{center}

\section{Верификация результатов}

Для проверки корректности результатов использовался Python и библиотека NumPy.

Результат программы сравнивается с результатом функции numpy.dot().

\section{Вывод}

В ходе выполнения лабораторной работы:

\begin{itemize}
\item реализована программа перемножения квадратных матриц
\item измерено время выполнения программы
\item выполнена автоматизированная проверка результатов
\item исследована зависимость времени выполнения от размера задачи
\end{itemize}

Алгоритм имеет вычислительную сложность $O(n^3)$.

\end{document}